\section*{第 \ref{ch:b2:isotope-geochemistry} 章——同位素地球化学}

\begin{solution}{exo:b2:isotope-geochemistry:1}
特里斯坦：$(0.51249/0.512638 - 1)\times10^4 = -2.9$。曼加亚：$(0.51285/0.512638 -
1)\times10^4 = +4.1$。MORB：$+8.5$。曼加亚与 MORB 的源区具有正的
$\varepsilon_{\mathrm{Nd}}$，平均而言其 Sm/Nd 高于\omterm{def:b2:planet-formation-meteorites:chondrule}{球粒陨石}；
特里斯坦的源区则较低。
\end{solution}

\begin{solution}{exo:b2:isotope-geochemistry:2}
$\delta^{18}\mathrm{O}_{\mathrm{VPDB}} = (\num{2.0672e-3}/\num{2.0052e-3} - 1)\times1000 =
+30.9$‰，以 \omterm{def:b2:isotope-geochemistry:standards}{VSMOW} 标度计。方解石的 $^{18}\mathrm{O}/^{16}\mathrm{O} =
\num{2.0672e-3}\times(1 - 0.0020) = \num{2.0631e-3}$，因此 $\delta^{18}\mathrm{O} =
(\num{2.0631e-3}/\num{2.0052e-3} - 1)\times1000 = +28.9$‰（以 \omterm{def:b2:isotope-geochemistry:standards}{VSMOW} 计）。
\end{solution}

\begin{solution}{exo:b2:isotope-geochemistry:3}
在 \qty{1.0}{Ga} 时：$e^{\lambda t} - 1 = 0.006545$，且 $0.512638 - 0.1966\times0.006545 =
0.511351$。在 \qty{4.567}{Ga} 时：$e^{\lambda t}-1 = 0.030243$，且 $0.512638 -
0.1966\times0.030243 = 0.506692$.
\end{solution}

\begin{solution}{exo:b2:isotope-geochemistry:4}
$\delta\mathrm{D} = 8\times(-18.0) + 10 = -134$‰。对湖水，$d = -42 -
8\times(-4.0) = -10$‰，低于直线上的 10：湖水位于\omterm{prop:b2:isotope-geochemistry:mwl}{大气降水
线}之下，与蒸发的地表水一致（轻同位素先逸出）。
\end{solution}

\begin{solution}{exo:b2:isotope-geochemistry:5}
$Q = 10^4\times\num{6.524e-12}\times0.1966/0.512638 = \num{2.502e-8}$ 每年，或
\qty{25.0}{\per\Ga}。花岗岩的 $\mathrm{Sm}/\mathrm{Nd} = 6.0/30 = 0.200$，
而 CI 为 $0.145/0.457 = 0.317$：$f = 0.200/0.317 - 1 = -0.37$。经过
\qty{0.5}{Ga}：$\Delta\varepsilon = 25.0\times(-0.37)\times0.5 = -4.6$ 个单位。
\end{solution}

\begin{solution}{exo:b2:isotope-geochemistry:6}
$\varepsilon \approx 3.5 + 25.0\times(-0.45)\times2.7 = 3.5 - 30.4 = -26.9$，而精确值为
$-27.3$：近似值偏高 0.4 个单位（负得不够多）。差异的大部分，
约 0.3 个单位，来自 $e^{\lambda t} > 1$：母体
过去比现在更丰富，因而相对于 CHUR 的生长亏损更大
（见\cref{exo:b2:isotope-geochemistry:10}）；其余部分来自 $f$ 的取整与初始值。
\end{solution}

\begin{solution}{exo:b2:isotope-geochemistry:7}
设花岗岩的质量分数为 $x$，则 $\mathrm{Sr} = 320x + 130(1-x)$，且
$R = (320x\times0.7200 + 130(1-x)\times0.7030)/\mathrm{Sr}$。当 $x = 0.05$ 时：Sr
为 \qty{139.5}{ppm}，$R = 0.70495$；$x = 0.10$：149.0，0.70665；$x = 0.20$：168.0，
0.70948。由混合命题的证明可知，$R\,\mathrm{Sr}$ 关于 Sr 是线性的：
$R = 0.731632 - 3.7221/\mathrm{Sr}$，这是一条斜率为 $-3.722$（ppm）的直线，
截距为 0.7316，横坐标为 $1/\mathrm{Sr}$。在 $x = 0.05$ 处检验：$0.731632 - 3.7221/139.5
= 0.70495$.
\end{solution}

\begin{solution}{exo:b2:isotope-geochemistry:8}
在 \qty{20}{\degreeCelsius} 时：$\delta^{18}\mathrm{O} = 1000/1.0098 - 1000 = -9.70$‰，
$\delta\mathrm{D} = 1000/1.084 - 1000 = -77.5$‰，斜率为 $77.5/9.70 = 7.99$。
在 \qty{0}{\degreeCelsius} 时：$-11.56$ 和 $-99.9$‰，斜率为 8.64。两者都接近
\omterm{prop:b2:isotope-geochemistry:mwl}{大气降水线}的斜率 8：水在液体与蒸气之间的平衡分馏
单独就决定了氘与氧-18 偏移量之比
接近 8。
\end{solution}

\begin{solution}{exo:b2:isotope-geochemistry:9}
地壳含有 $0.0051\times27/1.25 = 0.110$，即钕的 11\%。质量
平衡：$0.110\times(-15) + 0.890\,\varepsilon_m = 0$，故 $\varepsilon_m = 1.65/0.890 =
+1.9$。若亏损储库占有份额 $x$，其值为 $+8.5$，其余部分为 0：
$0.110\times(-15) + 8.5x = 0$，$x = 1.65/8.5 = 0.19$：约占
钕的五分之一。
\end{solution}

\begin{solution}{exo:b2:isotope-geochemistry:10}
生长速率为 $\lambda p\,e^{\lambda t}$，其中 $e^{\lambda t} \approx 1 + \lambda t$；
在从分离时刻 $T$ 到现在的区间上取平均，该因子为
$1 + \lambda T/2$，因此被忽略项的相对大小为 $\lambda T/2$。在
地球年龄范围内，$\num{6.524e-12}\times\num{4.567e9}/2 = 0.015$，即 1.5\%。对于
\cref{exo:b2:isotope-geochemistry:6} 中的地壳，$\lambda T/2 = 0.0088$，误差为 $0.0088\times30.4 = 0.27$
个单位。用来除 $\varepsilon$ 的 CHUR 比值的缓慢变化还带来
另一项较小的项，量级为 $(p/R)\lambda T \times \varepsilon_0$，此处为 0.02 个单位。
\end{solution}

\begin{solution}{exo:b2:isotope-geochemistry:11}
$K = 13.3/5.0 = 2.66$。取 $s = 0.5$：$s' = 2.66\times0.5/(1 + 1.66\times0.5) =
1.33/1.83 = 0.73$。差值 $s' - s$ 最大处满足
$\mathrm{d}s'/\mathrm{d}s = K/(1+(K-1)s)^2 = 1$，即 $s = (\sqrt K - 1)/(K-1) =
0.631/1.66 = 0.38$，此时 $s' = 0.62$：相差 0.24。
\end{solution}

\begin{solution}{exo:b2:isotope-geochemistry:12}
生成量：$2.3$ 至 $2.6\times\qty{7.2e9}{km^3}/\qty{4.5e9}{\text{年}} = 3.7$ 至
\qty{4.2}{km^3} 每年。消亡量：$1.3$ 至 $1.6\times7.2/4.5 = 2.1$ 至
\qty{2.6}{km^3} 每年。现今的速率有所不同：增加量（2.7）低于模型的
平均生成量，消减量（3.2）高于模型的平均消亡量。在模型中，地壳
早期增长，那时生成最快；而现在是平衡或略有亏损的时期，
这与俯冲带的收支结论相符。
\end{solution}

\begin{solution}{pb:b2:isotope-geochemistry:1}
\textbf{1.} 地幔：$(0.51307/0.512638 - 1)\times10^4 = +8.4$。沉积物：
$(0.51141/0.512638 - 1)\times10^4 = -24.0$.

\textbf{2.} Kora $+2.8$，Vel $-4.3$，Asha $+4.7$。

\textbf{3.} 以 $^{87}\mathrm{Sr}/^{86}\mathrm{Sr}$ 为横轴、$\varepsilon_{\mathrm{Nd}}$ 为纵轴：
地幔（0.70282，$+8.4$），Kora（0.70410，$+2.8$），Vel（0.70618，$-4.3$），沉积物
（0.72000，$-24.0$），Asha（0.70290，$+4.7$）。只有 Vel
的 $\varepsilon_{\mathrm{Nd}}$ 为负，位于 CHUR 的富集一侧。

\textbf{4.} 位于阵列之下：在\omterm{def:b2:isotope-geochemistry:reservoirs}{亏损地幔}的锶比值处，该阵列
应有 $\varepsilon_{\mathrm{Nd}}$ 接近 $+8$；Asha 为 $+4.7$。它与
HIMU 型的圣赫勒拿群（曼加亚、图布艾、圣赫勒拿）相似。

\textbf{5.} 熔化改变母体–子体比值，但不会可测地改变
锶和钕的同位素比值（关于熔化的命题）；熔岩
较年轻，因此自喷发以来其比值未曾增长（第 6 题）。

\textbf{6.} $\Delta R = 0.12\,(e^{\lambda\cdot10^6} - 1) = 0.12\times\num{6.52e-6} =
\num{7.8e-7}$，或 $\num{7.8e-7}/0.512638\times10^4 = 0.015$ 个 ε 单位：远低于
岛屿之间的差异，也低于一次测量的精度。

\textbf{7.} $\mathrm{Nd} = 40f + 0.75(1-f)$，且
$R = \bigl(40f\times0.51141 + 0.75(1-f)\times0.51307\bigr)/\mathrm{Nd}$.

\textbf{8.} $f = 0.5$\%：Nd \qty{0.9463}{ppm}，$\varepsilon = +1.6$；1\%：
1.1425，$-2.9$；2\%：1.535，$-8.4$；5\%：2.7125，$-15.4$。

\textbf{9.} Sr 10.95，11.90，13.80 和 \qty{19.50}{ppm}；$^{87}\mathrm{Sr}/^{86}\mathrm{Sr} =
0.70439$，0.70571，0.70780 和 0.71163。

\textbf{10.} $R\,\mathrm{Sr} = 10\times0.70282 + f(200\times0.72000 - 10\times0.70282)$ 和
$\mathrm{Sr} = 10 + 190f$；消去 $f$，$R\,\mathrm{Sr} = 7.0282 +
(\mathrm{Sr}-10)\times136.9718/190$，故 $R = 0.720904 - 0.18084/\mathrm{Sr}$：一条直线，
斜率为 $-0.1808$（ppm），对 $1/\mathrm{Sr}$ 而言。取 $f = 1$\% 验算：
$0.720904 - 0.18084/11.90 = 0.70571$.

\textbf{11.} $K = (10/0.75)/(200/40) = 13.33/5.0 = 2.67$。地幔相对
富锶，沉积物富钕，因此钕比值向沉积物一侧
移动得比锶比值快得多：少量沉积物就能使 $\varepsilon_{\mathrm{Nd}}$ 大幅移动，
而 $^{87}\mathrm{Sr}/^{86}\mathrm{Sr}$ 移动不大。

\textbf{12.} 钕含量高 53 倍（$40/0.75$），锶高 20 倍。在 $f = 1$\% 时，
沉积物提供了 $0.40/1.1425 = 35$\% 的钕，使
$\varepsilon_{\mathrm{Nd}}$ 改变 $0.35\times(8.4 + 24.0) = 11.3$ 个单位，而只提供了 $2.0/11.9 =
17$\% 的锶。

\textbf{13.} $0.01\times40/1.1425 = 0.350$：35\%。

\textbf{14.} 由 $R\,(40f + 0.75(1-f)) = 40fR_S + 0.75(1-f)R_D$：
$f = 0.75(R - R_D)/\bigl(40(R_S - R) + 0.75(R - R_D)\bigr)$，对锶则以 200
和 10 分别代替 40 和 0.75。

\textbf{15.} $f = 0.75\times(-0.00029)/\bigl(40\times(-0.00137) + 0.75\times(-0.00029)\bigr)
= \num{-2.175e-4}/\num{-5.502e-2} = 0.0040$：0.40\%。

\textbf{16.} $f = 10\times0.00128/\bigl(200\times0.01590 + 10\times0.00128\bigr) =
0.0128/3.1928 = 0.0040$：0.40\%；两者一致。

\textbf{17.} 钕：$0.75\times(-0.00065)/\bigl(40\times(-0.00101) +
0.75\times(-0.00065)\bigr) = \num{-4.875e-4}/\num{-4.089e-2} = 0.0119$；锶：
$10\times0.00336/\bigl(200\times0.01382 + 10\times0.00336\bigr) = 0.0336/2.7976 = 0.0120$。
两者均给出 1.2\%。

\textbf{18.} $f = 0.75\times(-0.00019)/\bigl(40\times(-0.00147) + 0.75\times(-0.00019)\bigr)
= 0.0024$，由此预测 Sr \qty{10.46}{ppm} 和 $^{87}\mathrm{Sr}/^{86}\mathrm{Sr} =
0.70361$，而实测为 0.70290。Asha 不在地幔–沉积物混合曲线上：
其源区含另一组分。

\textbf{19.} 古老的再循环洋壳，属 HIMU 型：其钕同位素要求 Sm/Nd
略低于\omterm{def:b2:isotope-geochemistry:reservoirs}{亏损地幔}的值，其锶同位素要求低 Rb/Sr，并且其
铅同位素（此处未测）要求高 U/Pb，这些过程均历时 10 亿至 20 亿年。

\textbf{20.} $e^{\lambda_{147}\cdot2\times10^9} - 1 = 0.013133$。沉积物：$0.51141 -
0.110\times0.013133 = 0.509965$；CHUR：$0.512638 - 0.1966\times0.013133 = 0.510056$；
$\varepsilon_{\mathrm{Nd}}(2\,\mathrm{Ga}) = (0.509965/0.510056 - 1)\times10^4 = -1.8$.

\textbf{21.} $f = 0.110/0.1966 - 1 = -0.440$；由于俯冲，$\varepsilon$ 变化了
$25.0\times(-0.440)\times2.0 = -22.0$，故 $\varepsilon(2\,\mathrm{Ga}) \approx -24.0 + 22.0
= -2.0$，与精确值相差在 0.2 个单位以内。

\textbf{22.} $e^{\lambda_{87}\cdot2\times10^9} - 1 = 0.028338$;
$0.72000 - 0.50\times0.028338 = 0.70583$.

\textbf{23.} 俯冲时，沉积物的 $\varepsilon_{\mathrm{Nd}} = -1.8$，接近
当时仅 7 亿年的地壳的 $-4.7$，并远高于
古老地壳的强负值：这与年轻地壳的侵蚀相符，或许
还混入了一些源自地幔的物质。当时的 $^{87}\mathrm{Sr}/^{86}\mathrm{Sr}$
为 0.7058，也是年轻地壳岩石的值。

\textbf{24.} Vel 的源区（最富集的岛屿）含有
\textbf{\qty[text-series-to-math=true]{1.2}{\percent}} 的再循环沉积物。
\end{solution}
